\begin{abstract}
多智能体大语言模型（LLM）系统通过输出聚合形成集体判断，但此类评审团是否表现出人类委员会中早有记录的羊群效应病理，以及若存在、哪类机制能够修复，目前仍缺乏深入理解。回答这两个问题需要两件仪器：测量集体行为的手段与干预集体行为的手段。本文表明，此前开发的博弈论共识框架恰好提供了这两件仪器：其 AEIC（假设-证据-推论-结论）分解通过显式结论字段使羊群效应成为可测量的量，其三层 Stackelberg 机制恰好映射到三个干预位置——聚合时校准、讨论时调控、逐实例协议路由。我们将框架置于"诊断后干预"的验证范式下，以自动化学术评审这一共识质量本质上关乎行为健康的领域为试验场。\textbf{研究一}对 300 篇分层采样的 ICLR~2024 论文进行诊断，每篇论文由 5 个同质 LLM 评审智能体组成的评审团在三种聚合条件下评审。独立 LLM 评审表现出系统性拒稿偏置（明确接收层论文的接收率仅 14.4\%）；三轮非结构化讨论触发\emph{负向棘轮}：少数派立场以 96:0 的比例向拒稿方向翻转（二项检验 $p{=}2.5{\times}10^{-29}$；消灭率 78\%），评审团评分方差坍缩 51\%（Wilcoxon $p{=}3.9{\times}10^{-32}$），均分显著下移——而与官方决议的决策一致率从未改善（61.7\%$\to$59.7\%，不显著）。跨模型复现（DeepSeek-V3、GLM-4-Flash、DeepSeek-R1）揭示了\emph{方向相反的病态吸引子}：V3 与 R1 评审团向拒稿方向棘轮化，而 GLM 评审团向接收方向棘轮化且 94\% 的结论为骑墙表述；推理能力最强的 R1（决策一致率 67.3\%）反而遭受最大的多样性坍缩（76\%）。\textbf{研究二}将框架的三个层次分别作为针对单一病理的干预。静态校准层将决策一致率从 60\% 提升至 67\%（双臂 66.7\%/67.3\%，合并 McNemar $p{=}0.034$），明确接收层从 2\% 修复至 44\%；Stackelberg 优化器在无人工提示下自主学到恒正校准项（$\tau{=}+1.04$），证明校准是博弈均衡的内生属性，其简约软实现（隐含阈值 4.88，两臂完全一致）消除了硬阈值调参的选择噪声。该精度距人类程序委员会的委员会间一致性（约 74\%，NeurIPS 一致性实验）仅 7 个百分点。提示层面的行为调控——证据锚定的少数派坚持——基本未能撼动棘轮（消灭率 78\%$\to$71\%），但保持了决策质量；这一诚实的负结果表明话术干预无法克服由能力驱动的从众。逐论文协议路由以低 56\% 的成本匹配最优固定协议（64.7\%），将讨论降级为按需调用的资源。两项研究共同确立了 LLM 共识系统的"校准优先"原则，并表明行为健康指标——少数派存活率、方差保持、漂移——是比决策准确率敏感得多的评审团级评估终点。

\keywords{多智能体共识 \and Stackelberg博弈 \and LLM同行评审 \and 羊群效应 \and 谄媚 \and 校准聚合}
\end{abstract}
